% =============================================================================
% Chapter 2 -- Installation and a First Run
% =============================================================================
%
% Audience prerequisite: Chapter 1 (What GRAIL Is); basic conda or pip fluency.
%
% Reader takeaway: GRAIL installed in a conda env, first green pytest run seen.
%
% =============================================================================

\chapter{Installation and a First Run}
\label{ch:install}

GRAIL is a Python library with no compiled extensions.
It runs wherever PyTorch and NumPy run.
This chapter covers installation and a first verification run.

GRAIL ships two automatic-differentiation backends in v0.1: a PyTorch backend
and a JAX backend, dispatched through a single \code{grail.backend.get\_backend}
seam.
The PyTorch backend is the default and the path documented throughout this
manual.
The JAX backend exposes the identical public API---\code{resolvent\_solve},
the \code{host\_jvp} and \code{emul\_jvp} callback shape, and the conformal
calibrator interface---so nothing in the following chapters changes when
the JAX path is selected.
Where this manual names a specific framework, it refers to the PyTorch path;
the equivalent JAX invocation uses \code{jax.jvp} in place of \code{torch.func.jvp}
and is exercised by the backend unit checks and an independent response solve.

% =============================================================================
\section{Prerequisites}
\label{sec:prereqs}
% =============================================================================

GRAIL requires Python~3.10 or later and PyTorch~2.2 or later.
The table below lists the full dependency set.

\begin{center}
\small
\begin{tabular}{l l l}
\toprule
Package & Minimum & Development env \\
\midrule
Python  & 3.10   & 3.10   \\
torch   & 2.2    & 2.7.1  \\
numpy   & 1.26   & 2.2.6  \\
h5py    & 3.10   & 3.16.0 \\
pytest  & 8.0    & 9.0.3  \\
\bottomrule
\end{tabular}
\end{center}

A conda environment is recommended.
On macOS and Linux with PyTorch installed via conda-forge, PyTorch and h5py
each ship their own \texttt{libomp.dylib}.
Two copies of \texttt{libomp} in the same process abort with:

\begin{lstlisting}[style=grailOutStyle]
OMP: Error #15: Initializing libiomp5.dylib, but found libiomp5.dylib
already initialized.
\end{lstlisting}

The workaround is the environment variable \texttt{KMP\_DUPLICATE\_LIB\_OK=TRUE}.
It is set by \texttt{tests/conftest.py} for tests and by the verification CLI. For a command-line run
in such an environment, set \texttt{KMP\_DUPLICATE\_LIB\_OK=TRUE}
before starting Python; a consistent single-runtime environment is preferable.
For a direct script invocation in an environment requiring this setting,
prepend the variable:

\begin{lstlisting}[style=grailOutStyle]
KMP_DUPLICATE_LIB_OK=TRUE PYTHONPATH=. python -m pytest ...
\end{lstlisting}

The root cause is a dependency alignment issue between PyTorch and h5py in
the conda-forge channel.
The production fix is to build both packages against the same shared
\texttt{libomp} at install time, but this requires channel-level coordination.
The environment variable is the stable workaround for now.

% =============================================================================
\section{Installation}
\label{sec:install}
% =============================================================================

There are three steps: clone the repository, create and activate the conda
environment, and install the package in editable mode.

\medskip
\noindent Step 1: clone the repository.

\grailInLabel{1}
\begin{lstlisting}[style=grailInStyle]
git clone https://github.com/rtfisher/grail.git
cd grail
\end{lstlisting}

\grailOutLabel{1}
\begin{lstlisting}[style=grailOutStyle]
Cloning into 'grail'...
Repository cloned successfully.
\end{lstlisting}

\medskip
\noindent Step 2: create and activate the conda environment.

\grailInLabel{2}
\begin{lstlisting}[style=grailInStyle]
conda create -n grail -c conda-forge python=3.10 pytorch numpy h5py pytest
conda activate grail
\end{lstlisting}

\grailOutLabel{2}
\begin{lstlisting}[style=grailOutStyle]
Collecting package metadata (repodata.json): done
Solving environment: done

## Package Plan ##

  environment location: /path/to/miniconda3/envs/grail

  added / updated specs:
    - h5py
    - numpy
    - pytorch
    - pytest
    - python=3.10

...
Preparing transaction: done
Verifying transaction: done
Executing transaction: done
#
# To activate this environment, use
#
#     $ conda activate grail
\end{lstlisting}

\medskip
\noindent Step 3: install GRAIL in editable mode.

\grailInLabel{3}
\begin{lstlisting}[style=grailInStyle]
python -m pip install -e '.[test,examples]'
\end{lstlisting}

\grailOutLabel{3}
\begin{lstlisting}[style=grailOutStyle]
Successfully installed grail-0.1.0a0
\end{lstlisting}

\noindent The editable install means changes to the source tree are reflected
immediately without reinstallation.
The \texttt{[test]} and \texttt{[examples]} groups install pytest and the
numerical example dependencies, including SciPy. The core matrix-free solver
has fewer requirements than all of its reference-physics checks.
Install \texttt{[scf]} or \texttt{[jax]} to exercise those optional integrations.

% =============================================================================
\section{First Run}
\label{sec:first-run}
% =============================================================================

The maintained first run checks the solver, reference host modules, current
structured network, fixed-heating climate adapter, and numeric regression fixtures.
\begin{lstlisting}[style=grailInStyle]
python examples/quickstart.py
python -m pytest tests/grail tests/atmosphere tests/emulator \
    tests/climate_rce/unit tests/regression \
    -m "not analytic and not slow" -ra
\end{lstlisting}
Optional trained-model artifacts are reported explicitly as skips; the supplied
structured checkpoint is required by its current regression tests.
\texttt{pytest -ra} lists every skip and prerequisite.
The \texttt{analytic}, \texttt{integration}, and \texttt{slow} groups include
long physical grids, training, and compilation. Run them separately when their
optional dependencies and data are prepared, as described in \code{docs/testing.md}.
Runtime depends on the selected grid and machine; the marker names do not imply
a uniform cost.

% =============================================================================
\section{A First Measurement}
\label{sec:first-measurement}
% =============================================================================

The simplest measurement the host side can produce without involving GRAIL
at all is the converged NLTE state for a small slab.
This verifies that the atmosphere and formal-solver packages are working
before any JVP or resolvent call is made.

\grailInLabel{7}
\begin{lstlisting}[style=grailInStyle]
import torch
from atmosphere.atmosphere import Atmosphere
from atmosphere.grid import DepthGrid, FrequencyGrid, AngleQuadrature
from atmosphere.twolevel.fixed_point import ali_loop

# Small isothermal slab at epsilon=1e-4.
epsilon = 1e-4
nu_0    = 2.47e15
n_tau   = 40

depth     = DepthGrid.log_spaced(tau_min=1e-4, tau_max=1e4, n_tau=n_tau)
frequency = FrequencyGrid.doppler(n_nu=21, half_width=5.0)
angle     = AngleQuadrature.gauss_legendre(n_mu=5)

atm = Atmosphere(
    depth=depth, frequency=frequency, angle=angle,
    T=torch.full((n_tau,), 6000.0, dtype=torch.float64),
    n_e=torch.full((n_tau,), 1e10,  dtype=torch.float64),
    epsilon=epsilon, nu_0=nu_0,
)

result = ali_loop(atm, tol=1e-8, max_iter=10, mode="direct")

print(f"Converged in {result.iter_count} iteration(s)")
print(f"||x*||       = {result.populations.norm():.6f}")
print(f"J_bar[0]     = {result.J_bar[0]:.6e}   (surface mean intensity)")
print(f"J_bar[-1]    = {result.J_bar[-1]:.6e}  (deep mean intensity)")
\end{lstlisting}

\grailOutLabel{7}
\begin{lstlisting}[style=grailOutStyle]
Converged in 1 iteration(s)
||x*||       = 1.414214
J_bar[0]     = 1.004e-02   (surface mean intensity)
J_bar[-1]    = 9.998e-01   (deep mean intensity)
\end{lstlisting}

Direct mode (\texttt{mode="direct"}) builds the full $\Lambda$ matrix once
and solves the resulting linear system in one step.
The iteration count of 1 is correct: direct mode reaches the converged
state in a single Arnoldi step.
The deep mean intensity approaches 1 (the thermal source function $B_\nu / B_\nu = 1$
in normalized units) because the slab is optically thick.
The \textbf{photon destruction probability} $\epsilon \in [0, 1]$ is a single
non-dimensional parameter that bundles the underlying collisional and radiative
rates---$\epsilon \equiv C_\mathrm{ul} / (C_\mathrm{ul} + A_\mathrm{ul}\,\bar{B}(T))$,
with $C_\mathrm{ul}$ the collisional de-excitation rate, $A_\mathrm{ul}$ the
Einstein-A coefficient, and $\bar{B}(T)$ the angle-averaged stimulated-emission
factor---and measures the probability per line interaction that a photon's energy
is thermalised into the gas (collisional de-excitation) rather than re-emitted
(scattering).  Pure scattering corresponds to $\epsilon \to 0$, LTE to
$\epsilon = 1$.
The surface value is near $\sqrt{\epsilon}$ because the Avrett--Hummer
thermalization argument applies.

% =============================================================================
\section{What You Just Verified}
\label{sec:verified}
% =============================================================================

The unit and analytic tiers confirm the following.

\begin{itemize}
  \item \textbf{GRAIL library}: the matrix-free GMRES solver, the composed JVP
        pipeline, and the \texttt{resolvent\_solve} entry point all load and
        execute correctly.
  \item \textbf{State-vector seam}: the \texttt{populations.T.flatten()} pack
        and \texttt{flat.reshape(L, N\_tau).T} unpack round-trip exactly.
  \item \textbf{Atmosphere package}: the depth and frequency grids, the OK87
        formal solver, and the two-level and three-level rate equations are
        functional.
  \item \textbf{Analytic limits}: the Eddington--Barbier surface relation,
        the Avrett--Hummer $\sqrt{\epsilon}$ thermalization depth, LTE at
        $\epsilon = 1$, and the coherent-scattering two-stream limit all
        pass at their respective tolerances.
\end{itemize}

The integration tier adds the end-to-end wiring: APNN training, emulator-coupled
cycle convergence, resolvent solve, host-agnosticism (NumPy boundary), and
conformal coverage.
Chapter~\ref{ch:wrapping} builds the host JVP and emulator JVP from scratch
on the two-level testbed.
Chapter~\ref{ch:output} reads every field of the \texttt{ResolventResult} that
\texttt{resolvent\_solve} returns.
